\section{Glueing properties}
\label{section-more-glueing}

\noindent
In this section we put a number of standard results of the
form: if something is true for all members of a standard open
covering then it is true. In fact, it often suffices to check
things on the level of local rings as in the following lemma.

\begin{lemma}
\label{lemma-characterize-zero-local}
Let $R$ be a ring.
\begin{enumerate}
\item For an element $x$ of an $R$-module $M$ the following are equivalent
\begin{enumerate}
\item $x = 0$,
\item $x$ maps to zero in $M_\mathfrak p$ for all $\mathfrak p \in \Spec(R)$,
\item $x$ maps to zero in $M_{\mathfrak m}$ for all maximal ideals
$\mathfrak m$ of $R$.
\end{enumerate}
In other words, the map $M \to \prod_{\mathfrak m} M_{\mathfrak m}$
is injective.
\item Given an $R$-module $M$ the following are equivalent
\begin{enumerate}
\item $M$ is zero,
\item $M_{\mathfrak p}$ is zero for all $\mathfrak p \in \Spec(R)$,
\item $M_{\mathfrak m}$ is zero for all maximal ideals $\mathfrak m$ of $R$.
\end{enumerate}
\item Given a complex $M_1 \to M_2 \to M_3$
of $R$-modules the following are equivalent
\begin{enumerate}
\item $M_1 \to M_2 \to M_3$ is exact,
\item for every prime $\mathfrak p$ of $R$ the localization
$M_{1, \mathfrak p} \to M_{2, \mathfrak p} \to M_{3, \mathfrak p}$
is exact,
\item for every maximal ideal $\mathfrak m$ of $R$ the localization
$M_{1, \mathfrak m} \to M_{2, \mathfrak m} \to M_{3, \mathfrak m}$
is exact.
\end{enumerate}
\item Given a map $f : M \to M'$ of $R$-modules the following are equivalent
\begin{enumerate}
\item $f$ is injective,
\item $f_{\mathfrak p} : M_\mathfrak p \to M'_\mathfrak p$ is injective
for all primes $\mathfrak p$ of $R$,
\item $f_{\mathfrak m} : M_\mathfrak m \to M'_\mathfrak m$ is injective
for all maximal ideals $\mathfrak m$ of $R$.
\end{enumerate}
\item Given a map $f : M \to M'$ of $R$-modules the following are equivalent
\begin{enumerate}
\item $f$ is surjective,
\item $f_{\mathfrak p} : M_\mathfrak p \to M'_\mathfrak p$ is surjective
for all primes $\mathfrak p$ of $R$,
\item $f_{\mathfrak m} : M_\mathfrak m \to M'_\mathfrak m$ is surjective
for all maximal ideals $\mathfrak m$ of $R$.
\end{enumerate}
\item Given a map $f : M \to M'$ of $R$-modules the following are equivalent
\begin{enumerate}
\item $f$ is bijective,
\item $f_{\mathfrak p} : M_\mathfrak p \to M'_\mathfrak p$ is bijective
for all primes $\mathfrak p$ of $R$,
\item $f_{\mathfrak m} : M_\mathfrak m \to M'_\mathfrak m$ is bijective
for all maximal ideals $\mathfrak m$ of $R$.
\end{enumerate}
\end{enumerate}
\end{lemma}

\begin{proof}
Let $x \in M$ as in (1). Let $I = \{f \in R \mid fx = 0\}$.
It is easy to see that $I$ is an ideal (it is the
annihilator of $x$). Condition (1)(c) means that for
all maximal ideals $\mathfrak m$ there exists an
$f \in R \setminus \mathfrak m$ such that $fx =0$.
In other words, $V(I)$ does not contain a closed point.
By Lemma \ref{lemma-Zariski-topology} we see $I$ is the unit ideal.
Hence $x$ is zero, i.e., (1)(a) holds. This proves (1).

\medskip\noindent
Part (2) follows by applying (1) to all elements of $M$ simultaneously.

\medskip\noindent
Proof of (3). Let $H$ be the homology of the sequence, i.e.,
$H = \Ker(M_2 \to M_3)/\Im(M_1 \to M_2)$. By
Proposition \ref{proposition-localization-exact}
we have that $H_\mathfrak p$ is the homology of the sequence
$M_{1, \mathfrak p} \to M_{2, \mathfrak p} \to M_{3, \mathfrak p}$.
Hence (3) is a consequence of (2).

\medskip\noindent
Parts (4) and (5) are special cases of (3). Part (6) follows
formally on combining (4) and (5).
\end{proof}

\begin{lemma}
\label{lemma-cover}
\begin{slogan}
Zariski-local properties of modules and algebras
\end{slogan}
Let $R$ be a ring. Let $M$ be an $R$-module. Let $S$ be an $R$-algebra.
Suppose that $f_1, \ldots, f_n$ is a finite list of
elements of $R$ such that $\bigcup D(f_i) = \Spec(R)$,
in other words $(f_1, \ldots, f_n) = R$.
\begin{enumerate}
\item If each $M_{f_i} = 0$ then $M = 0$.
\item If each $M_{f_i}$ is a finite $R_{f_i}$-module,
then $M$ is a finite $R$-module.
\item If each $M_{f_i}$ is a finitely presented $R_{f_i}$-module,
then $M$ is a finitely presented $R$-module.
\item Let $M \to N$ be a map of $R$-modules. If $M_{f_i} \to N_{f_i}$
is an isomorphism for each $i$ then $M \to N$ is an isomorphism.
\item Let $0 \to M'' \to M \to M' \to 0$ be a complex of $R$-modules.
If $0 \to M''_{f_i} \to M_{f_i} \to M'_{f_i} \to 0$ is exact for each $i$,
then $0 \to M'' \to M \to M' \to 0$ is exact.
\item If each $R_{f_i}$ is Noetherian, then $R$ is Noetherian.
\item If each $S_{f_i}$ is a finite type $R_{f_i}$-algebra, then
$S$ is a finite type $R$-algebra.
\item If each $S_{f_i}$ is of finite presentation over $R_{f_i}$, then
$S$ is a finitely presented $R$-algebra.
\end{enumerate}
\end{lemma}

\begin{proof}
We prove each of the parts in turn.
\begin{enumerate}
\item By Proposition \ref{proposition-localize-twice}
this implies $M_\mathfrak p = 0$ for all $\mathfrak p \in \Spec(R)$,
so we conclude by Lemma \ref{lemma-characterize-zero-local}.
\item For each $i$ take a finite generating set $X_i$ of $M_{f_i}$.
Without loss of generality, we may assume that the elements of $X_i$ are
in the image of the localization map $M \rightarrow M_{f_i}$, so we take
a finite set $Y_i$ of preimages of the elements of $X_i$ in $M$. Let $Y$
be the union of these sets. This is still a finite set.
Consider the obvious $R$-linear map $R^Y \rightarrow M$ sending the basis
element $e_y$ to $y$. By assumption this map is surjective after localizing
at an arbitrary prime ideal $\mathfrak p$ of $R$, so it is surjective by
Lemma \ref{lemma-characterize-zero-local}
and $M$ is finitely generated.
\item By (2) we have a short exact sequence
$$
0 \rightarrow K \rightarrow R^n \rightarrow M \rightarrow 0
$$
Since localization is an exact functor and $M_{f_i}$ is finitely
presented we see that $K_{f_i}$ is finitely generated for all
$1 \leq i \leq n$ by Lemma \ref{lemma-extension}.
By (2) this implies that $K$ is a finite $R$-module and therefore
$M$ is finitely presented.
\item By Proposition \ref{proposition-localize-twice}
the assumption implies that the induced morphism
on localizations at all prime ideals is an isomorphism, so we conclude
by Lemma \ref{lemma-characterize-zero-local}.
\item By Proposition \ref{proposition-localize-twice} the assumption
implies that the induced
sequence of localizations at all prime ideals is short exact, so we
conclude by Lemma \ref{lemma-characterize-zero-local}.
\item We will show that every ideal of $R$ has a finite generating set:
For this, let $I \subset R$ be an arbitrary ideal. By
Proposition \ref{proposition-localization-exact}
each $I_{f_i} \subset R_{f_i}$ is an ideal. These are all
finitely generated by assumption, so we conclude by (2).
\item For each $i$ take a finite generating set $X_i$ of $S_{f_i}$.
Without loss of generality, we may assume that the elements of $X_i$
are in the image of the localization map $S \rightarrow S_{f_i}$, so
we take a finite set $Y_i$ of preimages of the elements of $X_i$ in
$S$. Let $Y$ be the union of these sets. This is still a finite set.
Consider the algebra homomorphism $R[X_y]_{y \in Y} \rightarrow S$
induced by $Y$. Since it is an algebra homomorphism, the image $T$
is an $R$-submodule of the $R$-module $S$, so we can consider the
quotient module $S/T$. By assumption, this is zero if we localize
at the $f_i$, so it is zero by (1) and therefore $S$ is an
$R$-algebra of finite type.
\item By the previous item, there exists a surjective $R$-algebra
homomorphism $R[X_1, \ldots, X_n] \rightarrow S$. Let $K$ be the kernel
of this map. This is an ideal in $R[X_1, \ldots, X_n]$, finitely generated
in each localization at $f_i$. Since the $f_i$ generate the unit ideal
in $R$, they also generate the unit ideal in $R[X_1, \ldots, X_n]$, so an
application of (2) finishes the proof.
\end{enumerate}
\end{proof}

\begin{lemma}
\label{lemma-cover-upstairs}
Let $R \to S$ be a ring map.
Suppose that $g_1, \ldots, g_n$ is a finite list of
elements of $S$ such that $\bigcup D(g_i) = \Spec(S)$
in other words $(g_1, \ldots, g_n) = S$.
\begin{enumerate}
\item If each $S_{g_i}$ is of finite type over $R$, then $S$ is
of finite type over $R$.
\item If each $S_{g_i}$ is of finite presentation over $R$,
then $S$ is of finite presentation over $R$.
\end{enumerate}
\end{lemma}

\begin{proof}
Choose $h_1, \ldots, h_n \in S$ such that $\sum h_i g_i = 1$.

\medskip\noindent
Proof of (1). For each $i$ choose a finite list of elements
$x_{i, j} \in S_{g_i}$, $j = 1, \ldots, m_i$
which generate $S_{g_i}$ as an $R$-algebra.
Write $x_{i, j} = y_{i, j}/g_i^{n_{i, j}}$ for some $y_{i, j} \in S$ and
some $n_{i, j} \ge 0$. Consider the $R$-subalgebra $S' \subset S$
generated by $g_1, \ldots, g_n$, $h_1, \ldots, h_n$ and
$y_{i, j}$, $i = 1, \ldots, n$, $j = 1, \ldots, m_i$.
Since localization is exact (Proposition \ref{proposition-localization-exact}),
we see that $S'_{g_i} \to S_{g_i}$ is injective.
On the other hand, it is surjective by our choice of $y_{i, j}$.
The elements $g_1, \ldots, g_n$ generate the unit ideal in $S'$
as $h_1, \ldots, h_n \in S'$.
Thus $S' \to S$ viewed as an $S'$-module map is an isomorphism
by Lemma \ref{lemma-cover}.

\medskip\noindent
Proof of (2). We already know that $S$ is of finite type.
Write $S = R[x_1, \ldots, x_m]/J$ for some ideal $J$.
For each $i$ choose a lift $g'_i \in R[x_1, \ldots, x_m]$ of $g_i$
and we choose a lift $h'_i \in R[x_1, \ldots, x_m]$ of $h_i$.
Then we see that
$$
S_{g_i} = R[x_1, \ldots, x_m, y_i]/(J_i + (1 - y_ig'_i))
$$
where $J_i$ is the ideal of $R[x_1, \ldots, x_m, y_i]$
generated by $J$. Small detail omitted. By
Lemma \ref{lemma-finite-presentation-independent}
we may choose a finite list of elements
$f_{i, j} \in J$, $j = 1, \ldots, m_i$
such that the images of $f_{i, j}$ in $J_i$ and $1 - y_ig'_i$
generate the ideal $J_i + (1 - y_ig'_i)$.
Set
$$
S' = R[x_1, \ldots, x_m]/\left(\sum h'_ig'_i - 1, f_{i, j}; 
i = 1, \ldots, n, j = 1, \ldots, m_i\right)
$$
There is a surjective $R$-algebra map $S' \to S$.
The classes of the elements $g'_1, \ldots, g'_n$ in $S'$
generate the unit ideal and by construction the maps
$S'_{g'_i} \to S_{g_i}$ are injective.
Thus we conclude as in part (1).
\end{proof}





